\documentclass[11pt]{article}
\usepackage[T1]{fontenc}
\usepackage{lmodern,amsmath,amssymb,amsthm,bm,geometry,booktabs,longtable,hyperref}
\geometry{margin=24mm}
\hypersetup{colorlinks=true,linkcolor=blue,urlcolor=blue,citecolor=blue}
\setlength{\parskip}{0.35em}
\setlength{\emergencystretch}{3em}
\allowdisplaybreaks[2]
\newcommand{\NS}{\mathrm{NS}}
\newcommand{\R}{\mathrm{R}}
\newcommand{\FF}{\mathbb F}
\newcommand{\hFF}{\widehat{\mathbb F}}
\newcommand{\st}{\star_{\R}}
\newcommand{\norm}[1]{\left\|#1\right\|^2}
\numberwithin{equation}{section}
\title{Current algorithms for genus-two Ramond blocks\\
\large Ordinary theta sewing and recovery with a fermion insertion}
\author{Implementation notes}
\date{September 14, 2026}
\begin{document}
\maketitle

These notes describe the two production algorithms in \path{C++/}:
the ordinary NS--R--R theta block for $\eta\eta'=1$, and the
$\Theta v_{1/2}$ construction for $\eta\eta'=-1$.
The output in both cases is the physical superconformal block
$\FF_f^{(\eta,\eta')}$, including its dependence on the three spin lifts.
The shared ingredients are branching Ward identities, forward
central-charge recursion for the two ordinary Virasoro factors, the
Schottky vacuum, direct fermion sewing, and sector-restricted convolution.
Only the current algorithms and their implemented optimizations are described.
The mathematical statements concern generic parameters; numerical solves retain
rank and residual checks. No arbitrary-genus frontend is claimed for this code.

The computation proceeds in five steps: construct and reuse the finite action
data; evaluate and store the outer branching recurrence and, for the insertion,
the middle recurrence; evaluate the reduced Virasoro factors at branch-dependent cutoffs;
assemble the enlarged coefficients and divide by the direct fermion factor
in the appropriate parity sector; restore the common Schottky vacuum squared.
For the inserted graph only equal \emph{total} levels on its two middle
segments are needed. Section~\ref{sec:truncation} proves why this restriction
closes under recovery and gives the exact smaller domains needed inside the
two Virasoro factors.

\tableofcontents

\section{Conventions and the two recovery problems}
\label{sec:conventions}

The two spheres have punctures $(\infty,1,0)$ with sectors $(\NS,\R,\R)$.
Edges $1,2,3$ have plumbing parameters $q_1,q_2,q_3$ and spin lifts
$\eta_1,\eta_2,\eta_3\in\{\pm1\}$.
The vertex signs $\eta,\eta'$ are distinct from these spin lifts;
$f\in\{0,1\}$ labels the three-point form, and $p_1$ is the intrinsic
parity of the NS primary. We use the conventions of
\texttt{Human Notes/SCblock.tex}, with the contour transport specified below.
All contractions are bilinear BPZ contractions, without complex conjugating
the vertex coefficients. Hermitian norms appear only in numerical conditioning.
Total level $N$ means retaining $q_1^{a/2}q_2^{l_2}q_3^{l_3}$ with
$a,l_2,l_3\ge0$ integers and $a+2l_2+2l_3\le2N$.
The alternative independent-edge cutoff retains
$0\le a\le2N$ and $0\le l_2,l_3\le N$ without a bound on their sum.
The command-line choices are \texttt{--truncation total} (the default) and
\texttt{--truncation per-edge}; the latter applies the same numerical limit
$N$ independently to all three edges. The fresh timings below use this
second choice at $N=5$ and $N=10$, including the respective corners
$q_1^5q_2^5q_3^5$ and $q_1^{10}q_2^{10}q_3^{10}$.

Write
\begin{equation}
 Q=b+b^{-1},\qquad c=\frac32+3Q^2,\qquad
 h_1=\frac{Q^2/4-P_1^2}{2},\qquad
 \beta_j=P_j/\sqrt2\quad(j=2,3).
\end{equation}
The physical Ramond weights are $c/24-\beta_j^2$.
The implementation takes real $b\ne0,\pm1$ and permits complex momenta.
Genericity excludes singular embedding norms, unresolved Virasoro pole
collisions, and points where neither branching Ward identity determines the
required coefficient. A vanishing prefactor in the first identity alone is
handled by the second identity. Genuine singular limits require a separate
meromorphic-continuation prescription.

The physical Ramond ground conventions and three-point initial data are
\begin{gather}
 G_0w^\pm=i\beta e^{\mp i\pi/4}w^\mp,\qquad
 (\langle w^\alpha,w^\gamma\rangle)=\operatorname{diag}(1,i),\\
 \begin{array}{c|rrrr}
  & ++ & -- & +- & -+\\ \hline
  \rho_0^{(\eta)}(\phi,w^{\alpha},w^{\gamma})&1&\eta&0&0\\
  \rho_1^{(\eta)}(\phi,w^{\alpha},w^{\gamma})&0&0&1&i\eta
 \end{array}.
\end{gather}
The auxiliary grounds satisfy
$\psi_0u^{\mathsf b}=u^{1-\mathsf b}/\sqrt2$ and have metric
$\operatorname{diag}(1,-1)$. Tensor products are ordered auxiliary first;
a physical odd operator therefore acquires the auxiliary state's parity.

For parity monomials, the sewing convolution is
\begin{equation}
 \prod_i\eta_i^{\epsilon_i}\st\prod_i\eta_i^{\delta_i}
 =(-1)^{\sum_{i<j}(\epsilon_i\delta_j+\delta_i\epsilon_j)}
 \prod_i\eta_i^{\epsilon_i+\delta_i}.
 \label{eq:star}
\end{equation}
Parity exponents are reduced modulo two; plumbing exponents add ordinarily.
The cocycle is the difference between the quadratic theta sign at total
parity and the two separate quadratic signs. It is symmetric and bilinear,
so this product is commutative and associative.

The physical block obeys, coefficient by coefficient,
\begin{equation}
 (\eta_2\eta_3)\st\FF_f^{(\eta,\eta')}
 =\eta\eta'\FF_f^{(\eta,\eta')}.
 \label{eq:sector}
\end{equation}
One way to derive this is to change both Ramond summation bases by
\begin{equation}
 \mathbb L_{-B}w^+\mapsto(-1)^{\#G(B)}e^{-i\pi/4}\mathbb L_{-B}w^-,
 \qquad
 \mathbb L_{-B}w^-\mapsto-(-1)^{\#G(B)}e^{i\pi/4}\mathbb L_{-B}w^+.
 \label{eq:physical-intertwiner}
\end{equation}
Here $\#G(B)$ counts supercurrent modes. This odd intertwiner preserves the
ground metric $\operatorname{diag}(1,i)$ and supercommutes with the SCA.
Applying the map componentwise to both Ramond arguments multiplies a physical
three-point form by $-i\eta(-1)^{\epsilon_2}$.
The graded tensor action supplies another $(-1)^{\epsilon_2}$ and hence
has factor $-i\eta$. In the sewn pair of vertices the componentwise parity
factors cancel, giving $-\eta\eta'$.
Flipping the two parity bits changes the theta sign by
$(-1)^{\epsilon_2+\epsilon_3+1}$. Thus the coefficients at the flipped
and original parities differ by $\eta\eta'(-1)^{\epsilon_2+\epsilon_3}$,
which is exactly \eqref{eq:sector} by \eqref{eq:star}.

Ordinary sewing factorizes as
$\hFF_f^{(\eta,\eta')}=\FF_f^{(\eta,\eta')}\st\FF_{\mathsf F}$.
The uninserted auxiliary block satisfies
$(\eta_2\eta_3)\st\FF_{\mathsf F}=\FF_{\mathsf F}$ and has constant
$1+\eta_2\eta_3$. Its convolution with the negative sector vanishes in one line:
\begin{equation}
 \FF_{\mathsf F}\st\FF_f^{(\eta,-\eta)}
 =\FF_{\mathsf F}\st(\eta_2\eta_3)\st\FF_f^{(\eta,-\eta)}
 =-\FF_{\mathsf F}\st\FF_f^{(\eta,-\eta)}=0.
 \label{eq:cancellation}
\end{equation}
This is cancellation of the two Ramond ground contributions, equivalently
an unsaturated auxiliary zero mode in the corresponding loop trace. It is
not a vanishing statement about the physical block.
For the positive sector, the same auxiliary constant acts as the scalar $2$.
For the negative sector we instead construct an auxiliary constant
$Q(1-\eta_2\eta_3)/\sqrt2$, which acts as $\sqrt2Q$.
Neither constant is a unit in the entire parity algebra.

\section{Branching conventions and the enlarged blocks}
\label{sec:branch-sums}

The two commuting Virasoro algebras are generated by
\begin{align}
 L_m^{(1)}&=\frac{b^{-1}L_m-(b^{-1}+2b)L_m^{\mathsf F}+U_m}{b^{-1}-b},&
 L_m^{(2)}&=\frac{bL_m-(b+2b^{-1})L_m^{\mathsf F}+U_m}{b-b^{-1}},\\
 L_m^{\mathsf F}&=\frac12\sum_r r:\psi_{m-r}\psi_r:
                  +\frac{1-2\nu}{16}\delta_{m0},&
 U_m&=\sum_r\psi_{m-r}G_r.
\end{align}
Here $r\in\mathbb Z+\nu$, with $\nu=1/2$ in NS and $\nu=0$ in R.
We use $\psi_r^\dagger=-\psi_{-r}$, $L_m^\dagger=L_{-m}$,
$G_r^\dagger=G_{-r}$. The weights and central charges are
\begin{align}
 h_{j,n}^{(1)}&=\frac{Q^2/4-(P_j+2nb)^2}{2(1-b^2)},&
 c^{(1)}&=1+\frac{3Q^2}{1-b^2},\\
 h_{j,n}^{(2)}&=\frac{Q^2/4-(P_j+2n/b)^2}{2(1-b^{-2})},&
 c^{(2)}&=1+\frac{3Q^2}{1-b^{-2}}.
 \label{eq:weights}
\end{align}
The branch labels are $n_1\in\tfrac12\mathbb Z$ and
$n_2,n_3\in\tfrac12\mathbb Z+\tfrac14$. Their levels above the
enlarged ground states are $2n_1^2$ and $2n_j^2-1/8$, respectively.
In particular the Ramond subtraction is $1/8$, not the $1/16$ shift
between an enlarged weight and the physical Ramond weight.

\subsection{Primary normalization and the vertex frame}
Put $\chi_r(P)=\psi_r-i\eta_r(P)$, where $\eta_r$ is the physical
free-field fermion. For $n>0$ the NS primary is
\begin{equation}
 v_n(P)=2^{-2n}\ell(Q+2P,4n)
 \chi_{-1/2}(P)\chi_{-3/2}(P)\cdots\chi_{-(4n-1)/2}(P)\phi(P),
 \qquad v_0=\phi.
\end{equation}
This is the raw operator order used in C++. Reversing the $2n$ fermions
to the order in the human notes multiplies the primary by
$(-1)^{n(2n-1)}$. This phase changes raw vertices and action coefficients
consistently and cancels between the two sewn vertices.
For Ramond $n>0$, put $M=2n-1/2$ and let $t$ be the residue of $M+1$
modulo two. The two primaries are
\begin{align}
 v_n^t&=\chi_0(P)\chi_{-1}(P)\cdots\chi_{-M}(P)(u^0\otimes w^+),\\
 v_n^{1-t}&=\chi_0(P)\chi_{-1}(P)\cdots\chi_{-M}(P)
                 \chi_0(-P)(u^0\otimes w^+),
 \qquad \chi_0(-P)=\psi_0+i\eta_0(P).
\end{align}
Negative labels use simultaneous $n\mapsto-n$, $P\mapsto-P$;
on physical Ramond grounds the intertwiner is $\operatorname{diag}(1,-1)$.
This fixes the raw normalization, including the two parity copies.

For positive integer $m$,
\begin{equation}
 \ell(x,m)=2^{(m\bmod2)/8}
 \prod_{\substack{r,s\ge0,\ r+s<m\\r+s\equiv m\ (2)}}(x+rb+s/b),
 \qquad \ell(x,0)=1.
\end{equation}
For $n>0$, the squared BPZ norms used in assembly are
\begin{align}
 \norm{v_n}&=(-1)^{2n}2^{-2n}\ell(2P,4n)\ell(Q+2P,4n),\\
 \norm{v_n^0}&=2^{2\lfloor M/2\rfloor+1}
                 \frac{\ell(2P,4n)}{\ell(Q+2P,4n)},&
 \norm{v_n^1}&=-2^{2\lceil M/2\rceil}
                 \frac{\ell(2P,4n)}{\ell(Q+2P,4n)}.
 \label{eq:norms}
\end{align}
Negative labels again use reflection. Using squared norms avoids numerical
choices of square roots for normalized branching coefficients.

The code uses physical ground coordinates
$(w^+,e^{3\pi i/4}w^-)$, whose metric is $(1,1)$, and the Ward sign
$\eta_{\rm native}=(-1)^f\eta$.
After transporting these coordinates, the implemented enlarged vertex is
the following tensor-product form. Let $A,B$ be the physical descendant
words in the first two slots and $\mathsf A,\mathsf B,\mathsf C$ the auxiliary
words; $\alpha\in\{+,-\}$ is the physical second-slot ground label,
with $|+|=0$ and $|-|=1$. As in the SCblock notes, in sign exponents
$A,B$ denote the numbers of physical $G$ modes and
$\mathsf A,\mathsf B,\mathsf C$ the numbers of auxiliary $\psi$ modes:
\begin{align}
 \hat\rho_f^{(\eta)}=
 &(-i)^{A\bmod2}
 (-1)^{f(A+B+|\alpha|)+A\mathsf A}\nonumber\\
 &\times(-1)^{(B+|\alpha|+p_1)(\mathsf C+\mathsf c)}
 \rho_{\mathsf F}\rho_f^{(\eta)}.
 \label{eq:vertex-frame}
\end{align}
The arguments are in the order $(\NS,\R,\R)$, with auxiliary third ground
$u^{\mathsf c}$. The physical forms have the BPZ metric above and auxiliary
ground values $\rho_{\mathsf F}(\mathbf1,u^0,u^0)=1$,
$\rho_{\mathsf F}(\mathbf1,u^1,u^1)=i$.
Here the unhatted physical form uses the human ground and contour conventions
specified above; the powers of $-i$ and the $f$-dependent sign implement its
transport, rather than redefine its ground data.
These are parts of the definition, not adjustable phases in the final sum:
the two $(-i)^{A\bmod2}$ factors cancel the physical NS contour sign.

\subsection{Ordinary theta block}
All Virasoro blocks $F$ below have leading coefficient one, with primary
plumbing powers displayed separately. The ordinary enlarged block is
\begin{align}
 \hFF_f^{(\eta,\eta')}
 &=\sum_{\substack{n_1\in\frac12\mathbb Z,\ n_2,n_3\in\frac12\mathbb Z+\frac14\\
                         \alpha,\gamma\in\{0,1\}\\
                         2n_1+\alpha+\gamma\equiv f\ (2)}}
 q_1^{2n_1^2}q_2^{2n_2^2-1/8}q_3^{2n_3^2-1/8}
 \eta_1^{2n_1+p_1}\eta_2^\alpha\eta_3^\gamma\nonumber\\
 &\quad\times(-1)^{2n_1+(2n_1+p_1)(\alpha+\gamma)+\alpha\gamma}
 \frac{\hat\rho_f^{(\eta)}(v_{1,n_1}^{p_1},v_{2,n_2}^\alpha,v_{3,n_3}^\gamma)
       \hat\rho_f^{(\eta')}(v_{1,n_1}^{p_1},v_{2,n_2}^\alpha,v_{3,n_3}^\gamma)}
 {\norm{v_{1,n_1}^{p_1}}\norm{v_{2,n_2}^\alpha}\norm{v_{3,n_3}^\gamma}}
 \nonumber\\
 &\quad\times\prod_{i=1}^2
 F(h_{1,n_1}^{(i)},h_{2,n_2}^{(i)},h_{3,n_3}^{(i)};c^{(i)}\mid q_1,q_2,q_3).
 \label{eq:ordinary-sum}
\end{align}
The label ranges are those following \eqref{eq:weights}.
The fraction is the product of the two normalized branching coefficients
$\mathbb B_f^{(\eta);p_1}\mathbb B_f^{(\eta');p_1}$.
For $\eta\eta'=1$ this sum, divided by the ordinary auxiliary factor in
its sector, determines the physical block.

\subsection{Insertion for the negative Ramond loop}
Split edge 2 by a sphere with coordinate $z_3$:
\begin{equation}
 (z_1-1)z_3=\hat q_{2,1},\qquad
 (z_2-1)/z_3=\hat q_{2,2},\qquad
 q_2=\hat q_{2,1}\hat q_{2,2}.
 \label{eq:plumbing}
\end{equation}
Choose transported spin frames so that
$\eta_2=\hat\eta_{2,1}\hat\eta_{2,2}$.
In the irreducible physical NS vacuum module,
$G_{-1/2}\mathbf1=0$, so $v_{1/2}(Q/2)=Q\psi_{-1/2}\mathbf1$.
At $z_3=1$ insert the ordered operator $Q\Theta\psi(1)$, where
\begin{equation}
 \Theta\boldsymbol\Psi_{-\mathsf B}u^{\mathsf b}
 =(-1)^{\mathsf B+\mathsf b+1}
    \boldsymbol\Psi_{-\mathsf B}u^{1-\mathsf b}.
 \label{eq:Theta}
\end{equation}
It obeys $\Theta^2=-1$, anticommutes with $\psi_r$ and the physical $G_r$,
and commutes with $L_m$ and both $L_m^{(i)}$.
Thus $\Theta\psi$ is even and acts only on the auxiliary factor.
In particular
\begin{equation}
 Q\Theta\psi_0\boldsymbol\Psi_{-\mathsf B}u^{\mathsf b}
 =\frac Q{\sqrt2}(-1)^{\mathsf b}
     \boldsymbol\Psi_{-\mathsf B}u^{\mathsf b}.
 \label{eq:diagonal-insertion}
\end{equation}
The extra sign between Ramond grounds removes the cancellation in
\eqref{eq:cancellation}.

\paragraph{Definition in the PBW basis.}
Use the PBW notation of the SCblock notes. The enlarged NS and Ramond
basis states are, respectively,
\begin{align}
 &\Psi_{-\mathsf A}\mathbf L_{-A}\phi_1,\qquad
 \boldsymbol\Psi_{-\mathsf B}u^{\mathsf b}
              \otimes\mathbb L_{-B}w_j^\alpha,
 \qquad j=2,3.
 \label{eq:inserted-pbw-bases}
\end{align}
Here $\mathbf L_{-A}$ and $\mathbb L_{-B}$ are ordered physical PBW
words in $L_{-m}$ and $G_{-r}$, with no $G_0$ in a Ramond word.
The auxiliary words $\Psi_{-\mathsf A}$ and
$\boldsymbol\Psi_{-\mathsf B}$ contain distinct negative modes,
half-integer in NS and integer in R, with no $\psi_0$.
The NS notation suppresses $\mathbf1_{\mathsf F}\otimes$.
All physical ground labels run over $\{+,-\}$ and auxiliary ground
labels over $\{0,1\}$. Absolute values of words, such as $|A|$ and
$|\mathsf A|$, denote their levels. The parity shorthand is
\begin{gather}
 \mathrm{sgn}(A)=p_1+A+\mathsf A,\qquad
 \mathrm{sgn}(B)=B+|\alpha|+\mathsf B+\mathsf b,\nonumber\\
 \mathrm{sgn}(C)=C+|\gamma|+\mathsf C+\mathsf c
 \pmod2.
 \label{eq:inserted-pbw-grading}
\end{gather}
The physical inverse Gram matrices are
$\mathbf G_{h_1}^{AA'}$ and
$\mathbb G_{\beta_j}^{B\alpha,B'\alpha'}$.
The auxiliary Gram matrices are diagonal:
\begin{align}
 \langle\Psi_{-\mathsf A}\mathbf1_{\mathsf F}
   |\Psi_{-\mathsf A'}\mathbf1_{\mathsf F}\rangle
 &=(-1)^{\mathsf A}\delta_{\mathsf A,\mathsf A'},\nonumber\\
 \langle\boldsymbol\Psi_{-\mathsf B}u^{\mathsf b}
   |\boldsymbol\Psi_{-\mathsf B'}u^{\mathsf b'}\rangle
 &=(-1)^{\mathsf B+\mathsf b}
       \delta_{\mathsf B,\mathsf B'}\delta_{\mathsf b,\mathsf b'}.
 \label{eq:inserted-pbw-grams}
\end{align}
The Kronecker symbols equate complete mode words. Tensor-product BPZ
pairings are products of the two pairings, without an extra graded sign.
We therefore sum out the primed auxiliary labels below and display their
inverse norms as signs.

Keep $A,A'$ on edge 1 and $C,C'$ on edge 3. On the two halves of
edge 2 use $B,B'$ and $D,D'$, with auxiliary words $\mathsf B,\mathsf D$
and grounds $(\alpha,\mathsf b),(\delta,\mathsf d)$, respectively;
$\mathrm{sgn}(D)=D+|\delta|+\mathsf D+\mathsf d\pmod2$.
With fixed ground-state plumbing powers suppressed, the full inserted
enlarged block is
\begin{align}
 &\hFF_f^{(\eta,\eta')}[\Theta v_{1/2}]
 =\sum
 q_1^{|A|+|\mathsf A|}
 \hat q_{2,1}^{|B|+|\mathsf B|}
 \hat q_{2,2}^{|D|+|\mathsf D|}
 q_3^{|C|+|\mathsf C|}\nonumber\\
 &\quad\times\eta_1^{\mathrm{sgn}(A)}
 \hat\eta_{2,1}^{\mathrm{sgn}(B)}
 \hat\eta_{2,2}^{\mathrm{sgn}(D)}
 \eta_3^{\mathrm{sgn}(C)}
 (-1)^{A+\mathsf A}
 (-1)^{\mathsf A+\mathsf B+\mathsf b+\mathsf D+\mathsf d+\mathsf C+\mathsf c}
 \nonumber\\
 &\quad\times
 (-1)^{\mathrm{sgn}(A)\mathrm{sgn}(B)+
        \mathrm{sgn}(A)\mathrm{sgn}(C)+\mathrm{sgn}(B)\mathrm{sgn}(C)}
 \mathbf G_{h_1}^{AA'}\mathbb G_{\beta_2}^{B\alpha,B'\alpha'}
 \mathbb G_{\beta_2}^{D\delta,D'\delta'}
 \mathbb G_{\beta_3}^{C\gamma,C'\gamma'}\nonumber\\
 &\quad\times\hat\rho_f^{(\eta)}\bigl(
 \Psi_{-\mathsf A}\mathbf L_{-A}\phi_1,
 \boldsymbol\Psi_{-\mathsf B}u^{\mathsf b}\otimes\mathbb L_{-B}w_2^\alpha,
 \boldsymbol\Psi_{-\mathsf C}u^{\mathsf c}\otimes\mathbb L_{-C}w_3^\gamma\bigr)
 \nonumber\\
 &\quad\times\hat\rho_f^{(\eta')}\bigl(
 \Psi_{-\mathsf A}\mathbf L_{-A'}\phi_1,
 \boldsymbol\Psi_{-\mathsf D}u^{\mathsf d}\otimes\mathbb L_{-D}w_2^\delta,
 \boldsymbol\Psi_{-\mathsf C}u^{\mathsf c}\otimes\mathbb L_{-C'}w_3^{\gamma'}\bigr)
 \nonumber\\
 &\quad\times(\mathbb G_{\beta_2})_{D'\delta',B'\alpha'}
 \langle\boldsymbol\Psi_{-\mathsf D}u^{\mathsf d}|
 Q\Theta\psi(1)|\boldsymbol\Psi_{-\mathsf B}u^{\mathsf b}\rangle.
 \label{eq:inserted-pbw-definition}
\end{align}
The sum runs over all displayed descendant and ground labels, with
$|A|=|A'|$, $|B|=|B'|$, $|D|=|D'|$, $|C|=|C'|$, and
\begin{gather}
 A+B+C+|\alpha|+|\gamma|\equiv f,\qquad
 A'+D+C'+|\delta|+|\gamma'|\equiv f,\nonumber\\
 \mathsf A+\mathsf B+\mathsf C+\mathsf b+\mathsf c\equiv0,\qquad
 \mathsf A+\mathsf D+\mathsf C+\mathsf d+\mathsf c\equiv0\pmod2.
 \label{eq:inserted-pbw-parities}
\end{gather}
Both outer forms use the transported convention \eqref{eq:vertex-frame}.
The signs on the second line of \eqref{eq:inserted-pbw-definition} are
the NS contour factor and the four auxiliary inverse norms; the next line
contains the theta sewing sign. The middle bra and ket are at
$z_3=\infty$ and $z_3=0$. Since the insertion acts only on the auxiliary
factor, their physical pairing is the Gram matrix with lower indices
displayed on the last line. In particular,
\begin{equation}
 \sum_{B',\alpha',D',\delta'}
 \mathbb G_{\beta_2}^{B\alpha,B'\alpha'}
 (\mathbb G_{\beta_2})_{D'\delta',B'\alpha'}
 \mathbb G_{\beta_2}^{D\delta,D'\delta'}
 =\mathbb G_{\beta_2}^{B\alpha,D\delta}.
 \label{eq:inserted-pbw-collapse}
\end{equation}
Thus the two halves have equal physical levels and parities.
The even insertion also gives $\mathrm{sgn}(B)=\mathrm{sgn}(D)$, so the
two spin lifts combine into $\eta_2^{\mathrm{sgn}(B)}$.
There is no restriction on their total levels in the full block:
$\psi(1)=\sum_{r\in\mathbb Z}\psi_r$ can change the auxiliary level.

For the diagonal used in production, retain
$|B|+|\mathsf B|=|D|+|\mathsf D|$ and replace the two middle plumbing
powers by $q_2^{|B|+|\mathsf B|}$. Equal physical levels then imply
equal auxiliary levels, selecting only $\psi_0$. Equation~\eqref{eq:diagonal-insertion}
gives
\begin{equation}
 \langle\boldsymbol\Psi_{-\mathsf D}u^{\mathsf d}|
 Q\Theta\psi_0|\boldsymbol\Psi_{-\mathsf B}u^{\mathsf b}\rangle
 =\frac Q{\sqrt2}(-1)^{\mathsf b}(-1)^{\mathsf B+\mathsf b}
 \delta_{\mathsf D,\mathsf B}\delta_{\mathsf d,\mathsf b}.
 \label{eq:inserted-pbw-middle}
\end{equation}
The surviving auxiliary inverse norms multiply to one by the even
auxiliary parity condition. Contracting the physical middle Gram matrix
and relabeling $D,\delta$ as $B',\alpha'$ therefore gives
\begin{align}
 &\operatorname{Diag}\hFF_f^{(\eta,\eta')}[\Theta v_{1/2}]
 =\sum
 q_1^{|A|+|\mathsf A|}q_2^{|B|+|\mathsf B|}q_3^{|C|+|\mathsf C|}
 \eta_1^{\mathrm{sgn}(A)}\eta_2^{\mathrm{sgn}(B)}\eta_3^{\mathrm{sgn}(C)}
 \nonumber\\
 &\quad\times
 \frac Q{\sqrt2}(-1)^{\mathsf b}
 (-1)^{\mathrm{sgn}(A)\mathrm{sgn}(B)+\mathrm{sgn}(A)\mathrm{sgn}(C)+
                  \mathrm{sgn}(B)\mathrm{sgn}(C)+A+\mathsf A}
 \mathbf G_{h_1}^{AA'}\mathbb G_{\beta_2}^{B\alpha,B'\alpha'}
 \mathbb G_{\beta_3}^{C\gamma,C'\gamma'}\nonumber\\
 &\quad\times\hat\rho_f^{(\eta)}\bigl(
 \Psi_{-\mathsf A}\mathbf L_{-A}\phi_1,
 \boldsymbol\Psi_{-\mathsf B}u^{\mathsf b}\otimes\mathbb L_{-B}w_2^\alpha,
 \boldsymbol\Psi_{-\mathsf C}u^{\mathsf c}\otimes\mathbb L_{-C}w_3^\gamma\bigr)
 \nonumber\\
 &\quad\times\hat\rho_f^{(\eta')}\bigl(
 \Psi_{-\mathsf A}\mathbf L_{-A'}\phi_1,
 \boldsymbol\Psi_{-\mathsf B}u^{\mathsf b}\otimes\mathbb L_{-B'}w_2^{\alpha'},
 \boldsymbol\Psi_{-\mathsf C}u^{\mathsf c}\otimes\mathbb L_{-C'}w_3^{\gamma'}\bigr).
 \label{eq:inserted-pbw-diagonal}
\end{align}
Here the sum has $|A|=|A'|$, $|B|=|B'|$, $|C|=|C'|$ and the parity
conditions \eqref{eq:inserted-pbw-parities} with
$D,\delta,\mathsf D,\mathsf d$ replaced by
$B',\alpha',\mathsf B,\mathsf b$.
Omitting $Q(-1)^{\mathsf b}/\sqrt2$ gives the ordinary enlarged theta
PBW sum of the SCblock notes, with the vertex transport
\eqref{eq:vertex-frame} understood.
The code evaluates the enlarged coefficients using the double-Virasoro
sum below. Substituting
\eqref{eq:vertex-frame} into \eqref{eq:inserted-pbw-diagonal} gives the
auxiliary factor $Q(-1)^{\mathsf b}/\sqrt2$ in
\eqref{eq:fermion-sum} and the convolution
\eqref{eq:diagonal-convolution}. Replacing $Q\Theta\psi(1)$ by the
identity in \eqref{eq:inserted-pbw-definition} instead collapses both
halves to the ordinary enlarged theta block at
$q_2=\hat q_{2,1}\hat q_{2,2}$.

\paragraph{Double-Virasoro decomposition.}
The inserted double-Virasoro field has weights
\begin{equation}
 h_\psi^{(1)}=-\frac{1+2b^2}{2(1-b^2)},\qquad
 h_\psi^{(2)}=\frac{b^2+2}{2(1-b^2)}.
\end{equation}
These have level-two degenerate Kac labels $(2,1)$ in the first copy and
$(1,2)$ in the second, in the embedding convention. At common physical momentum
their two fusion rules restrict the middle branches to
$n'=n\pm1/2$. Both halves carry the same enlarged parity $\alpha$,
because the field and $\Theta$ each reverse it.
Let the punctured Virasoro block have ordered internal weights
$(h_1,h_L,h_R,h_3)$ and external weight $h_\psi$; its outer vertices are
$(h_1,h_L,h_3)$ and $(h_1,h_R,h_3)$, and its middle vertex is
$(h_R,h_\psi,h_L)$ at $(\infty,1,0)$.
The enlarged block is
\begin{align}
 &\hFF_f^{(\eta,\eta')}[\Theta v_{1/2}]
 =\sum_{\substack{n_1\in\frac12\mathbb Z,\ n,n',n_3\in\frac12\mathbb Z+\frac14\\
 \alpha,\gamma\in\{0,1\}\\
 n'=n\pm1/2,\ 2n_1+\alpha+\gamma\equiv f\ (2)}}
 q_1^{2n_1^2}\hat q_{2,1}^{2n^2-1/8}
 \hat q_{2,2}^{2n'^2-1/8}q_3^{2n_3^2-1/8}
 \eta_1^{2n_1+p_1}\eta_2^\alpha\eta_3^\gamma\nonumber\\
 &\quad\times(-1)^{2n_1+(2n_1+p_1)(\alpha+\gamma)+\alpha\gamma}
 \frac{\hat\rho_f^{(\eta)}(v_{1,n_1}^{p_1},v_{2,n}^\alpha,v_{3,n_3}^\gamma)
       \hat\rho_f^{(\eta')}(v_{1,n_1}^{p_1},v_{2,n'}^\alpha,v_{3,n_3}^\gamma)}
 {\norm{v_{1,n_1}^{p_1}}\norm{v_{2,n}^\alpha}\norm{v_{2,n'}^\alpha}
  \norm{v_{3,n_3}^\gamma}}\nonumber\\
 &\quad\times\langle v_{2,n'}^\alpha|Q\Theta\psi(1)|v_{2,n}^\alpha\rangle
 \prod_{i=1}^2
 F(h_{1,n_1}^{(i)},h_{2,n}^{(i)},h_{2,n'}^{(i)},h_{3,n_3}^{(i)};
   h_\psi^{(i)};c^{(i)}\mid q_1,\hat q_{2,1},\hat q_{2,2},q_3).
 \label{eq:inserted-sum}
\end{align}
The additional inverse primary norm belongs to the additional cut.
The numerical calculation evaluates only the shifted diagonal of this
sewing expression, on the domains in Section~\ref{sec:truncation}.
The middle matrix element is a branching coefficient: in normalized radial
ordering it is $i\|v_{1/2}\|
\mathbb B(n',1/2,n;\alpha,1-\alpha)$ times the two internal norm roots,
with compatible roots satisfying $\Theta(v_n^\alpha/\|v_n^\alpha\|)
=-i v_n^{1-\alpha}/\|v_n^{1-\alpha}\|$.
Equation \eqref{eq:inserted-sum} uses raw matrix elements and is independent
of that optional square-root notation. Here $\norm{v_{1/2}(Q/2)}=-Q^2$.

\section{Computing and reusing the branching data}
\label{sec:branching}

\subsection{Finite action systems}
The columns are constructed from oscillators with
\begin{equation}
 [c_m,c_n]=m\delta_{m+n,0},\qquad
 \{\eta_r,\eta_s\}=\{\psi_r,\psi_s\}=\delta_{r+s,0},\qquad
 \{\eta_r,\psi_s\}=0.
\end{equation}
On the positive chart the physical free-field realization is
\begin{align}
 L_m&=\frac12\sum_{k\ne0,m}c_kc_{m-k}
       +\frac12\sum_r r\eta_{m-r}\eta_r
       +\frac i2(Qm-2P)c_m\quad(m\ne0),\\
 G_r&=\sum_{k\ne0}c_k\eta_{r-k}+i(Qr-P)\eta_r,\\
 L_0&=\frac{Q^2/4-P^2}{2}+\frac{1-2\nu}{16}
       +\sum_{k>0}c_{-k}c_k+\sum_{r>0}r\eta_{-r}\eta_r.
\end{align}
Positive modes annihilate the ground states. On the internal Ramond physical
coordinates, $\eta_0$ swaps the grounds with coefficient $1/\sqrt2$.
Reflection replaces $P$ by $-P$ and transports the ground as specified above.
Only finitely many terms of each mode sum act nontrivially on a given state.

For positive labels the required action expansions are
\begin{align}
 L_1v_n&=\sum_{|A|+|B|=4n-3}\mathbf V_{1,n}^{AB}
                 L_{-A}^{(1)}L_{-B}^{(2)}v_{n-1}\quad(n\ge1),\\
 L_1v_n^\alpha&=\sum_{|A|+|B|=4n-3}\mathbb V_{1,n,\alpha}^{AB}
                 L_{-A}^{(1)}L_{-B}^{(2)}v_{n-1}^\alpha\quad(n\ge3/4),\\
 L_{-1}v_n^\alpha&=\sum_{i=1}^2\mathbb V_{n,\alpha}^{(i)}L_{-1}^{(i)}v_n^\alpha
 +\sum_{|A|+|B|=4n-1}\mathbb V_{-1,n,\alpha}^{AB}
                 L_{-A}^{(1)}L_{-B}^{(2)}v_{n-1}^\alpha\quad(n\ge3/4).
 \label{eq:actions}
\end{align}
At NS $n=0,\pm1/2$ and R $n=\pm1/4$, $L_1$ annihilates the primary.
At $n=1/4$, the last line instead has the two same-primary level-one terms
and a primary term $v_{-3/4}^\alpha$; reflection supplies $n=-1/4$.
The code constructs this small boundary span directly. It does not extend
the generic $n\mapsto n-1$ formula across this boundary.

Expand the target on the left and every candidate descendant on the right
in the free-field oscillator basis. Solve the resulting finite column-span
problem for $\mathbf V$ or $\mathbb V$. The candidate supports in the
displayed action expansions are the structural input from the branching identities;
the implemented all-row residual checks test that the computed target lies
in those supports at the requested parameters. They are not a proof of
these identities at arbitrary level.
The ordinary pipeline requires the NS $L_1$ and Ramond $L_{-1}$ data.
The insertion additionally requires Ramond $L_1$ data on edge 2.
When the second outer Ward identity is needed, NS $L_{-1}$ and Ramond
$L_0,L_1$ actions are constructed lazily and stored. Here physical $L_0$
acts with the physical oscillator level, including the Ramond ground weight;
it is not replaced by the enlarged conformal weight.

\subsection{Outer branching from directly computed boundary values}
\label{sec:outer-recursion}
First directly compute and store all allowed coefficients in the boundary box
\begin{equation}
 n_1=0,\pm\tfrac12,\qquad n_2,n_3=\pm\tfrac14,\pm\tfrac34.
\end{equation}
This is the boundary prescription in \texttt{Human Notes/SCblock.tex}.
The explicit free-field states are converted to small physical PBW states;
their physical and auxiliary three-point forms follow from the Ward identities
with \eqref{eq:vertex-frame}. Only these initial values and low boundary
conversions use physical PBW. No high-level physical Gram sewing enters.

For higher labels, use the physical conformal Ward identities
\begin{align}
 \hat\rho(L_1v_{1,n_1},v_{2,n_2},v_{3,n_3})
 &=\hat\rho(v_{1,n_1},L_{-1}v_{2,n_2},v_{3,n_3})
  +\hat\rho(v_{1,n_1},v_{2,n_2},L_{-1}v_{3,n_3}),
 \label{eq:outer-ward}\\
 \hat\rho(v_{1,n_1},L_1v_{2,n_2},v_{3,n_3})
 &=\hat\rho(L_{-1}v_{1,n_1},v_{2,n_2},v_{3,n_3})
   -\hat\rho(v_{1,n_1},L_{-1}v_{2,n_2},v_{3,n_3})\nonumber\\
 &\quad-2\hat\rho(v_{1,n_1},L_0v_{2,n_2},v_{3,n_3})
   -\hat\rho(v_{1,n_1},v_{2,n_2},L_1v_{3,n_3}).
 \label{eq:outer-ward-second}
\end{align}
Insert the mode-action decompositions and factor descendant forms into the
raw primary coefficient and the two ordinary Virasoro three-point factors.
The first identity usually determines a coefficient from stored lower-label
values by a scalar division. If its prefactor vanishes, use the second
identity. Store every computed coefficient immediately and reuse the same
Ward coefficients for both requested vertex signs.

The actual dependency order increases
\begin{equation}
 \lfloor |n_1|\rfloor+\lfloor |n_2|-\tfrac14\rfloor
                         +\lfloor |n_3|-\tfrac14\rfloor.
\end{equation}
The boundary action exchanges the pairs $\{1/4,-3/4\}$ and
$\{-1/4,3/4\}$ on a Ramond leg; the second identity may also exchange
$\{1/2,-1/2\}$ on the NS leg. These exchanges preserve the displayed order.
Group such coupled coefficients at fixed order and substitute every stored
lower-order value first. Outside the directly seeded box, at most four
unknowns remain in a group. Solve this local relation at the working precision,
then continue upward. Thus the bulk update is scalar and the small boundary
couplings never become a global Ward matrix. Every external dependency is
checked to have smaller order and to be present in the cache.

Pivot tests use the contributions to the local prefactor and the other
unknowns in its small group. Large already-known lower-label contributions
are not included in this scaling. If the first identities do not determine
the group, append the corresponding second identities and retry locally.
After each update, evaluate all used identities with the stored coefficients.
An unresolved pivot or failed residual reports the particular labels, parities,
and, for a residual failure, vertex sign.

Start from the primary triples required by the branch sum and close this set
under the action-label changes. Retain
$2n_1+\alpha+\gamma\equiv f\pmod2$. For the inserted graph, initial candidates
use the average of the two middle primary levels under a total cutoff; the
stricter maximum-level condition in Section~\ref{sec:truncation} is imposed
during diagonal assembly. Under an independent-edge cutoff, both middle
primary levels must already be at most $N$, and the other primary levels
are bounded on their own edges. The closure may add lower-label coefficients
needed by the Ward recurrence beyond these initial branch requests.
The action stage still prepares its edge-label ranges before this closure.
The ordinary pipeline requests one vertex sign; the inserted pipeline requests
both, with the same local elimination applied to their two right-hand sides.

For a word $A=(a_1,\ldots,a_k)$, put $s_j=\sum_{t>j}a_t$.
The normalized single-leg descendant factors are the closed products
\begin{align}
 \rho(L_{-A}\nu_1,\nu_2,\nu_3)&=\prod_j(h_1+s_j+a_jh_2-h_3),\\
 \rho(\nu_1,L_{-A}\nu_2,\nu_3)&=\prod_j(-1)^{a_j}(h_2+s_j+a_jh_3-h_1),\\
 \rho(\nu_1,\nu_2,L_{-A}\nu_3)&=\prod_j(h_3+s_j+a_jh_2-h_1).
 \label{eq:single-leg}
\end{align}
These replace recursive three-point evaluation in the outer and middle
Ward equations. Products with the same changed labels, slot, copy, and word
are cached in the outer recurrence. The embedded weights are cached by leg,
label, and Virasoro copy; empty descendant words return one immediately.

\subsection{Middle recurrence}
Since $[L_m,Q\Theta\psi(1)]=0$ for the physical stress tensor,
\begin{equation}
 \langle L_1v_{n'}^\alpha|Q\Theta\psi(1)|v_n^\alpha\rangle
 =\langle v_{n'}^\alpha|Q\Theta\psi(1)|L_{-1}v_n^\alpha\rangle,
 \label{eq:middle-ward}
\end{equation}
and the transposed identity also holds. For $n'=n+1/2$ with $n\ge1/4$,
substitution gives
\begin{align}
 &\left[\sum_{i=1}^2\mathbb V_{n,\alpha}^{(i)}
 (h_{2,n}^{(i)}+h_\psi^{(i)}-h_{2,n'}^{(i)})\right]
 \langle v_{n'}^\alpha|Q\Theta\psi(1)|v_n^\alpha\rangle\nonumber\\
 &=\sum_{|A|+|B|=4n'-3}\mathbb V_{1,n',\alpha}^{AB}
 \rho(L_{-A}\nu_{h_{2,n'-1}^{(1)}},\nu_{h_\psi^{(1)}},\nu_{h_{2,n}^{(1)}})
 \rho(L_{-B}\nu_{h_{2,n'-1}^{(2)}},\nu_{h_\psi^{(2)}},\nu_{h_{2,n}^{(2)}})
 \langle v_{n'-1}^\alpha|Q\Theta\psi(1)|v_n^\alpha\rangle.
 \label{eq:middle-recursion}
\end{align}
The other-primary term of $L_{-1}v_n$ is separated from $n'$ by $3/2$
and vanishes by fusion. For the reverse orientation use the transposed
identity, with the descendant on the incoming leg. More generally apply
$L_1$ to the endpoint of larger absolute label, using reflection for negative
labels. Each step decreases that absolute label and reaches an oriented
ground pair. Equation \eqref{eq:diagonal-insertion} fixes both orientations:
\begin{equation}
 \langle v_{1/4}^0|Q\Theta\psi|v_{-1/4}^0\rangle=\sqrt2Q,
 \qquad
 \langle v_{1/4}^1|Q\Theta\psi|v_{-1/4}^1\rangle=Q/\sqrt2.
\end{equation}
The reversed pairs have the same values. The recurrence uses cached actions
from the outer stage and stores each ordered pair and parity once.

\subsection{Implemented reuse and numerical linear algebra}
\label{sec:branch-optimizations}
The following optimizations are part of the current C++ implementation.
\begin{enumerate}
\item \textbf{Reflection and parity transport.} Negative-label actions are
computed on the positive chart at $-P$ and relabelled, avoiding a large
Fock-to-PBW reflection matrix. Only parity-zero Ramond action systems are solved.
For a term $n\to m$, its parity-one coefficient is multiplied by
\begin{equation}
 2^{\{(-1)^{2|m|-1/2}-(-1)^{2|n|-1/2}\}/2}.
\end{equation}
This follows by commuting the action through
$\Theta v_n^0=-2^{(-1)^{2|n|-1/2}/2}v_n^1$.
It includes the factors $1/2$ and $2$ at the $1/4\leftrightarrow-3/4$ boundary.
\item \textbf{Interned sparse states.} Each oscillator occupation pattern has
one integer identifier. Sparse linear combinations store identifiers and
coefficients. Nonzero fermion occupations use bit masks; their reordering signs
come from population counts. Physical-mode images ignore auxiliary spectators;
auxiliary stress-tensor images ignore physical spectators. Spectators and graded
signs are restored when the stored image is used.
The fixed oscillator storage checks its mode bounds explicitly
(64 positive fermion slots and boson modes $1,\ldots,63$).
\item \textbf{Commuting-copy scheduling.} Descendant words share stored suffix
states and embedded-generator images. At each step the code selects the larger
available leading mode from the two commuting Virasoro words. This changes
intermediate work, not either word's internal order. It uses commutation of the
two copies; it does not identify their different weights or numerical blocks.
\item \textbf{Paired actions and cache lifetimes.} Required $L_{-1}$ and $L_1$
actions on edge 2 share primary and descendant construction before their
temporary image/suffix caches are released. Solved actions, primaries, and
interned states remain available to the run. No precomputed numerical file is
loaded. The two physical momenta and their reflected charts remain distinct.
\item \textbf{Stored branching recurrence.} Each directly evaluated boundary
value and each subsequently computed coefficient is retained. Ward rows,
embedded weights, and single-leg descendant factors are reused. The two
vertex signs share the same local elimination; its dimension is at most four.
\item \textbf{Scaled action systems.} Columns are normalized once. Native
action solves use full-column-rank LAPACK least squares. Multiprecision action
systems use pivoted QR of the transposed scaled matrix to select independent
rows, then an LU factorization of that square restriction.
\item \textbf{Multiprecision refinement.} Double precision supplies row choices
and an LU preconditioner only. Residuals are evaluated using the original
multiprecision columns, correction equations use the reusable preconditioner,
and corrections are accumulated at working precision. The final residual is
checked on every original row, including unselected rows, using sparse entries.
This obtains accurate coefficients without factoring a large dense MPC matrix.
\end{enumerate}

\section{Forward CCY recursion and the Schottky vacuum}
\label{sec:ccy}

Both Virasoro factors use the central-charge recursion of Cho, Collier and
Yin~\cite{CCY}. Separate the universal vacuum factor $F_{\rm vac}$ and write
$F=F_{\rm vac}F_{\rm red}$. At descendant multi-level $\boldsymbol\ell$,
\begin{equation}
 F_{{\rm red},\boldsymbol\ell}(c;\boldsymbol h)
 =G_{\boldsymbol\ell}(\boldsymbol h)
 +\sum_e\sum_{\substack{r\ge2\\s\ge1}}
 \frac{R_{rs}^{(e)}(\boldsymbol h)}{c-c_{rs}(h_e)}
 F_{{\rm red},\boldsymbol\ell-rs\boldsymbol e_e}
       (c_{rs}(h_e);\boldsymbol h+rs\boldsymbol e_e).
 \label{eq:ccy}
\end{equation}
Only nonnegative remaining levels contribute; $\boldsymbol e_e$ is the unit
vector on internal edge $e$, and $G$ is the global $SL(2)$ coefficient.
The external inserted weights remain fixed during all recursion steps,
including steps that change the central charge.

\subsection{Residues and their reusable factors}
With $x=b_{rs}(h)^2$,
\begin{align}
 x&=\frac{rs-1+2h\pm\sqrt{(r-s)^2+4(rs-1)h+4h^2}}{1-r^2},&
 c_{rs}(h)&=13+6(x+x^{-1}),\\
 R_{rs}^{(e)}&=-\partial_{h_e}c_{rs}(h_e)A_{rs}
                       P_{rs}(h_a,h_b)P_{rs}(h_d,h_k),\\
 A_{rs}&=\frac12\prod_{\substack{u=1-r,\ldots,r;\ v=1-s,\ldots,s\\
                  (u,v)\ne(0,0),(r,s)}}(ub_{rs}+v/b_{rs})^{-1}.
 \label{eq:residue}
\end{align}
The two fusion polynomials use the spectator pairs at the two endpoints:
\begin{center}
\begin{tabular}{c cc}
\toprule
Null edge & Ordinary theta (both endpoints) & Inserted graph (two endpoints)\\
\midrule
1 & $(h_3,h_2)$ & $(h_3,h_L)$, $(h_3,h_R)$\\
2 or $L$ & $(h_1,h_3)$ & $(h_3,h_1)$, $(h_R,h_\psi)$\\
$R$ & --- & $(h_3,h_1)$, $(h_L,h_\psi)$\\
3 & $(h_1,h_2)$ & $(h_1,h_L)$, $(h_1,h_R)$\\
\bottomrule
\end{tabular}
\end{center}
The order of each pair specifies the sign at odd null level.
For $\lambda_a^2=x+2+x^{-1}-4h_a$ the fusion polynomial is
\begin{equation}
 P_{rs}(h_a,h_b)=
 \prod_{\substack{u=1-r,3-r,\ldots,r-1\\v=1-s,3-s,\ldots,s-1}}
 \frac{\lambda_a+\lambda_b+ub_{rs}+v/b_{rs}}2
 \frac{\lambda_a-\lambda_b+ub_{rs}+v/b_{rs}}2.
\end{equation}
The code pairs $(u,v)$ with $(-u,-v)$, evaluating their product as
\begin{equation}
 \frac{\bigl(4(h_b-h_a)+u^2x+2uv+v^2/x\bigr)^2
       -4\lambda_a^2(u^2x+2uv+v^2/x)}{16}.
\end{equation}
An unpaired $(0,0)$ contributes $h_b-h_a$. This removes square-root choices
from fusion evaluation. The larger-magnitude quadratic root for $x$ avoids
a cancelling root. For $r,s\ge2$ the two roots have product
$(1-s^2)/(1-r^2)$; swapping $r,s$ makes the selected root the reciprocal
of the other one. Thus $(r,s)$ and $(s,r)$ together supply both finite
central-charge poles. For $r=s$ the roots are reciprocal and give one pole.
For $s=1$ one root is zero, so the nonzero root supplies the unique finite
pole; there is no additional $r=1$ contribution.
The derivative-times-$A$ product is simplified algebraically before evaluation:
its factors $x-1$ and $x+1$ cancel matching factors in the denominator.
Unresolved zero denominators are errors, not skipped residues.

Pole geometry and the combined universal residue are cached by null edge,
integer weight shift, $r,s$. Fusion polynomials additionally key the spectator
weights and shifts, including their order. These caches last for one ordinary
Virasoro block; different outer branching tuples instantiate separate engines.

\subsection{The global coefficients}
For one normalized Virasoro three-point form, let
$\rho(i,j,k)=\rho(L_{-1}^i\nu_1,L_{-1}^j\nu_2,L_{-1}^k\nu_3)$.
With $(x)_m$ rising and $x^{\underline m}$ falling factorials,
\begin{align}
 \rho(i,j,k)&=(h_1-h_2-h_3+i-j+1-k)_j\nonumber\\
 &\quad\times\sum_{p=0}^{\min(i,k)}\binom{i}{p}
 (2h_3+k-1)^{\underline p}k^{\underline p}
 (h_3+h_2-h_1)_{k-p}
 (h_1+h_2-h_3-k+p)_{i-p}.
 \label{eq:global-rho}
\end{align}
The ordinary seed is $\rho(a,b,d)^2/[a!(2h_1)_a b!(2h_2)_b d!(2h_3)_d]$.
The inserted seed is
\begin{equation}
 G_{a,b,c,d}=
 \frac{\rho_{1L3}(a,b,d)\rho_{1R3}(a,c,d)\rho_{R\psi L}(c,0,b)}
 {a!(2h_1)_a b!(2h_L)_b c!(2h_R)_c d!(2h_3)_d}.
 \label{eq:global-seed}
\end{equation}
Each subscript gives the three ordered weights.
The finite sum in \eqref{eq:global-rho} is independent of $j$ and is cached
without that index. Rising/falling products are built incrementally rather
than by repeated products inside the sum. Each global norm is cached separately
and appears once per internal edge in \eqref{eq:global-seed}.

For the inserted seed, three further caches retain the \emph{complete}
vertex factors, including the $j$-dependent rising product and their assigned
norms. The first outer vertex includes the inverse norms of edges $1,3$;
the second outer vertex carries no inverse norms; the middle vertex includes those
of $L,R$. Their product is exactly \eqref{eq:global-seed}. Each cache is
keyed by its three incident weight shifts and three descendant levels;
the external weight is fixed, with zero recursion shift and descendant level.
Thus a change on a nonincident edge
reuses the entire factor, including its normalization. The cache returns
a reference to the stored value. For MPC arithmetic, the final product of
these three factors and the incoming amplitude uses one reusable temporary,
avoiding repeated allocation of multiprecision mantissas in the assembly sum.
All these numerical caches belong to one Virasoro engine and are released
when that engine finishes; different branching weights remain separate.

\subsection{One forward pass for all coefficients}
The transition in \eqref{eq:ccy} depends on accumulated internal shifts and
the current central charge, but not on the final requested coefficient.
Store the amplitude $w_{\boldsymbol s}(c_*)$ of all paths reaching shifted
weights $\boldsymbol h+\boldsymbol s$ at central charge $c_*$.
Initialize $w_{\boldsymbol0}(c)=1$. In increasing total shift, update
\begin{equation}
 w_{\boldsymbol s+rs\boldsymbol e_e}(c_{rs}(h_e+s_e))\mathrel{+}=
 R_{rs}^{(e)}(\boldsymbol h+\boldsymbol s)
 \sum_{c_*}\frac{w_{\boldsymbol s}(c_*)}{c_*-c_{rs}(h_e+s_e)}.
 \label{eq:forward}
\end{equation}
Keep $c_*$ separate during propagation: summing it away would lose the
different outgoing denominators. The common residue is multiplied once after
the incoming sum. The reciprocal of each denominator is computed once and
stored by the ordered pair of incoming and outgoing central-charge identifiers;
subsequent uses multiply by that reciprocal. The small-denominator guard is
applied when the reciprocal is first constructed. This preserves the exact
recurrence, although factoring out the residue changes floating-point rounding.
After outgoing transitions have been processed,
store $\sum_{c_*}w_{\boldsymbol s}(c_*)$ and release that shift's incoming map.
Finally,
\begin{equation}
 F_{{\rm red},\boldsymbol\ell}(c;\boldsymbol h)
 =\sum_{\boldsymbol s\le\boldsymbol\ell}
       \left(\sum_{c_*}w_{\boldsymbol s}(c_*)\right)
       G_{\boldsymbol\ell-\boldsymbol s}(\boldsymbol h+\boldsymbol s).
 \label{eq:forward-assemble}
\end{equation}
This is the finite expansion of \eqref{eq:ccy}, grouped by final shift and
terminal global coefficient. Every transition increases a shift by $rs\ge2$,
so a downward-closed finite level set makes the forward pass acyclic.
It reuses all internal recursion paths across requested coefficients;
it does not reuse a completed block at different external parameters.
An exactly zero sum of incoming amplitudes is not used to discard transitions
with different $c_*$. Exact zero residues may be skipped.

\subsection{Exact Schottky vacuum}
The universal large-$c$ vacuum is computed from CCY Eq.~(5.5):
\begin{equation}
 F_{\rm vac}=\prod_{\gamma\in\mathcal P}\prod_{m\ge2}
                    (1-q_\gamma^m)^{-1/2}.
 \label{eq:schottky}
\end{equation}
Here $\mathcal P$ includes primitive conjugacy classes and their inverses
separately. The code identifies inverse pairs and uses exponent $-1$.
For the theta graph the three crossing matrices are
\begin{equation}
 \begin{pmatrix}0&1\\q_1&0\end{pmatrix},\qquad
 \begin{pmatrix}1&q_2-1\\1&-1\end{pmatrix},\qquad
 \begin{pmatrix}0&q_3\\1&0\end{pmatrix}.
\end{equation}
Enumerate directed closed walks with no immediate backtracking, modulo cyclic
rotation and inversion, excluding proper powers. The direction records the
starting sphere, so an odd-length path cannot be mistaken for a closed root.
For the projective matrix $M_\gamma$ of a retained walk,
\begin{equation}
 q_\gamma=\sum_{j\ge1}\frac1{j+1}\binom{2j}{j}
     \left(\frac{\det M_\gamma}{(\operatorname{tr}M_\gamma)^2}\right)^j.
\end{equation}
The trace is a unit at the theta cusp and the leading determinant degree is
the walk length. Since the first oscillator is $m=2$, only walks of length
at most $\lfloor N/2\rfloor$ can contribute through total level $N$.
For the independent-edge cutoff, a walk traversing edge $i$ a total of
$v_i$ times can contribute only if $2v_i\le N$ on every edge, hence its
length is at most $\lfloor3N/2\rfloor$. After enumerating primitive classes
through this length, the code discards any class exceeding an edge-visit bound
before computing its multiplier. All multiplier expansions, products,
logarithms and exponentials are truncated componentwise in this case.
Multiply the factors using their logarithmic series and a truncated exponential,
entirely over GMP rationals. This is a cutoff derived from degree, not a
separate empirical limit on word length.
The inserted sphere collapses by M\"obius maps when its light external
weight is omitted from this universal factor, giving the same vacuum with
$q_2=\hat q_{2,1}\hat q_{2,2}$. Both reduced Virasoro factors omit one
vacuum, so two copies are restored after physical recovery.

\section{Exact target truncation and products}
\label{sec:truncation}

In the physical factor, the inserted field is the identity. Writing $G$ for
its Gram matrix, the middle contraction is $G^{-1}GG^{-1}=G^{-1}$ at each
level and parity.
Consequently the full factorization has the same physical theta block:
\begin{equation}
 \hFF_f^{(\eta,-\eta)}[\Theta v_{1/2}]
 =\FF_f^{(\eta,-\eta)}(q_1,\hat q_{2,1}\hat q_{2,2},q_3)
       \st\FF_{\mathsf F}[\Theta v_{1/2}].
 \label{eq:full-factorization}
\end{equation}
Let $\operatorname{Diag}$ retain equal powers of the two split parameters
and replace their product by $q_2$. Since the physical factor has equal powers,
\begin{equation}
 \operatorname{Diag}\hFF_f^{(\eta,-\eta)}[\Theta v_{1/2}]
 =\FF_f^{(\eta,-\eta)}\st
       \operatorname{Diag}\FF_{\mathsf F}[\Theta v_{1/2}].
 \label{eq:diagonal-convolution}
\end{equation}
Indeed, at a diagonal target $(k,k)$ the physical contribution $(j,j)$
forces the auxiliary contribution to be $(k-j,k-j)$. This proves closure;
no off-diagonal physical coefficient must first be computed or divided.
The auxiliary diagonal selects precisely the zero-mode action
\eqref{eq:diagonal-insertion}. Unequal middle levels in each Virasoro
factor in \eqref{eq:inserted-sum} are nevertheless required on the domain
below, before taking the shifted diagonal of their product.

\subsection{Total-level cutoff}
For a fixed inserted branch define the integer budgets
\begin{equation}
 K=\left\lfloor N-2n_1^2-(2n_3^2-1/8)\right\rfloor,\qquad
 L=K-(2n^2-1/8),\qquad R=K-(2n'^2-1/8).
\end{equation}
Discard branches with $\min(L,R)<0$. In each individual Virasoro factor
keep exactly the downward closure
\begin{equation}
 a+b+d\le L,\qquad a+c+d\le R,\qquad a,b,c,d\ge0.
 \label{eq:target-domain}
\end{equation}
Indeed, a lower descendant index extends to an allowed shifted diagonal
if and only if
$a+d+\max(2n^2-1/8+b,\,2n'^2-1/8+c)\le K$:
choose the common final middle power to be this maximum.
This criterion is exactly the two inequalities above.
For a product target put $t=0,\ldots,\min(L,R)$, choose $a+d\le t$,
and set $(b,c)=(L-t,R-t)$. These are all required targets. For each, sum
only the splits of that target between the two Virasoro series.
The closure is preserved by every CCY subtraction. Each factor needs unequal
middle levels inside it: imposing $b=c$ separately would discard valid terms.
The largest allowed middle level is $N$ in the lowest branch, not $2N$.

For ordinary theta sewing, each Virasoro factor instead uses the simplex
\begin{equation}
 a+b+d\le\left\lfloor N-2n_1^2-(2n_2^2-1/8)-(2n_3^2-1/8)\right\rfloor.
\end{equation}
Thus there is a different cutoff for every branching triple in either mode.

\subsection{Independent-edge cutoff}
For an inserted branch, put
\begin{equation}
 L=N-(2n^2-1/8),\qquad R=N-(2n'^2-1/8).
\end{equation}
Each individual Virasoro factor is evaluated at nonnegative integer
descendant levels in the box
\begin{equation}
 \begin{gathered}
 0\le a\le\lfloor N-2n_1^2\rfloor,\qquad
 0\le d\le N-(2n_3^2-1/8),\\
 0\le b\le L,\qquad 0\le c\le R.
 \end{gathered}
 \label{eq:target-box}
\end{equation}
Discard a branch if any upper bound is negative. In the product of the two
factors retain the same first- and third-edge bounds and impose
\begin{equation}
 2n^2-1/8+b=2n'^2-1/8+c\le N.
 \label{eq:box-diagonal}
\end{equation}
Here $a,b,c,d$ in the product are the sums of the two factors' descendant
indices. Equation~\eqref{eq:target-box} is precisely the downward closure
of these shifted-diagonal targets: any index in the box can be completed
by increasing its smaller middle total to the larger one, leaving its first
and third levels unchanged. Conversely every lower index of an allowed
target lies in the box. Thus individual factors need unequal middle levels,
but neither middle cutoff is doubled. Product assembly sums only the
splittings of the required targets.

For ordinary sewing the corresponding three-dimensional box is
\begin{equation}
 0\le a\le\lfloor N-2n_1^2\rfloor,\qquad
 0\le b\le N-(2n_2^2-1/8),\qquad
 0\le d\le N-(2n_3^2-1/8).
\end{equation}
Every stage uses these componentwise bounds, including branching requests,
CCY, products, auxiliary sewing, division and vacuum restoration. The code
does not first compute a total-level-$3N$ block and then discard coefficients.
The physical box has $(2N+1)(N+1)^2$ monomials: $2{,}541$ at $N=10$,
compared with $506$ at total level 10. Both truncations are downward closed,
so the forward recursion and triangular recovery close on their requested sets.

\subsection{Reuse under middle-edge transposition}
BPZ transposition at the insertion point one gives
\begin{align}
 &F(h_1,h_L,h_R,h_3;h_\psi;c\mid q_1,q_L,q_R,q_3)\nonumber\\
 &\qquad=F(h_1,h_R,h_L,h_3;h_\psi;c\mid q_1,q_R,q_L,q_3).
\end{align}
The code computes each product once for an unordered pair $\{n,n'\}$ at
fixed $n_1,n_3$, then transposes its middle exponents for the reverse orientation.
Both orientations retain their own branching coefficients, norms, and parity
factors. The stored product is released after its second use.

\section{Direct fermion factor and physical recovery}
\label{sec:fermion}

Use distinct descending positive half-integer modes $\mathsf A$ on the NS
edge and distinct descending positive integer modes $\mathsf B,\mathsf C$
on Ramond edges. Include $u^{\mathsf b},u^{\mathsf c}$ with
$\mathsf b,\mathsf c\in\{0,1\}$. Put
$\delta_1=\#\mathsf A$, $\delta_2=\#\mathsf B+\mathsf b$,
$\delta_3=\#\mathsf C+\mathsf c$ modulo two.
Only $\delta_1+\delta_2+\delta_3=0$ contributes. In the following state sum,
use $1$ for ordinary sewing and $Q(-1)^{\mathsf b}/\sqrt2$ for the inserted
diagonal:
\begin{align}
 &\sum_{\mathsf A,\mathsf B,\mathsf C,\mathsf b,\mathsf c}
 q_1^{|\mathsf A|}q_2^{|\mathsf B|}q_3^{|\mathsf C|}
 \eta_1^{\delta_1}\eta_2^{\delta_2}\eta_3^{\delta_3}
 (-1)^{\delta_1+\delta_1\delta_2+\delta_1\delta_3+\delta_2\delta_3}
 \nonumber\\[-2pt]
 &\quad\times
 \rho_{\mathsf F}(\boldsymbol\Psi_{-\mathsf A}\mathbf1,
                  \boldsymbol\Psi_{-\mathsf B}u^{\mathsf b},
                  \boldsymbol\Psi_{-\mathsf C}u^{\mathsf c})^2
 \begin{cases}1,&\text{ordinary},\\Q(-1)^{\mathsf b}/\sqrt2,&\text{inserted diagonal}.
 \end{cases}
 \label{eq:fermion-sum}
\end{align}
Each Fock norm is $(-1)^{\delta_i}$. In ordinary sewing their inverse
product is one on an allowed even-parity triple. In inserted diagonal sewing
the additional edge-2 inverse norm cancels the BPZ norm in the middle matrix
element, leaving $Q(-1)^{\mathsf b}/\sqrt2$.
For a total cutoff retain $|\mathsf A|+|\mathsf B|+|\mathsf C|\le N$;
for an independent-edge cutoff bound each of $|\mathsf A|,|\mathsf B|,
|\mathsf C|$ separately by $N$.
The constant terms are $1+\eta_2\eta_3$ and
$Q(1-\eta_2\eta_3)/\sqrt2$, respectively.

Three-point forms are evaluated by the fermion Ward identity. Remove the largest
mode in the first nonempty slot and move its contour using, for target NS
mode $k-1/2$ or Ramond mode $n$, respectively,
\begin{equation}
 z^{k-1/2}(z-1)^{-1/2},\qquad
 z^{-1/2}(z-1)^{-n-1/2},\qquad
 z^{-n-1/2}(z-1)^{-1/2}.
\end{equation}
The local expansions give half-integer binomial coefficients, frame phases
$(i,1,-i)$, and graded contour signs. Solve for the original form; other terms
have lower nonzero-mode level. Stop each local expansion when its modes
annihilate the states, with a bound computed from the largest occupied mode.
Store every resulting three-point form across the entire sewing sum.

For exact arithmetic temporarily rescale $u^1\to\sqrt2u^1$ and strip
the phase $i^{\delta_2}$. The stored rational form satisfies
\begin{equation}
 \rho_{\mathsf F}=i^{\delta_2}2^{-(\mathsf b+\mathsf c)/2}
                      \rho_{\mathsf F}^{\rm rat}.
\end{equation}
Its ground values are $1,2$, and its zero-mode matrix entries are $1/2,1$.
All Ward arithmetic and sewing are then rational. Restore only the common
$Q/\sqrt2$ for the inserted factor at the chosen numerical precision.
Only diagonal zero-mode transitions are enumerated in production.

For either sector, subtract positive-level auxiliary convolutions using already
solved lower coefficients. Divide the remainder by $2$ for ordinary sewing or
$\sqrt2Q$ for inserted diagonal sewing. Explicitly, for any physical monomial
$q^{\boldsymbol k}$ and the appropriate numerator/auxiliary pair,
\begin{equation}
 \FF[q^{\boldsymbol k}]=
 \frac{\hFF[q^{\boldsymbol k}]
 -\sum_{\substack{\boldsymbol0\le\boldsymbol r\le\boldsymbol k\\
                  \boldsymbol r\ne\boldsymbol0}}
       \FF_{\mathsf F}[q^{\boldsymbol r}]\st
       \FF[q^{\boldsymbol k-\boldsymbol r}]}{
       \begin{cases}2,&\eta\eta'=1,\\\sqrt2Q,&\eta\eta'=-1.\end{cases}}
 \label{eq:division}
\end{equation}
The numerator on the negative side means its diagonal projection. With the
reduced Virasoro sum as input, this recurrence first gives
$\FF/F_{\rm vac}^2$; multiply by $F_{\rm vac}^2$ afterwards.
The sector identity \eqref{eq:sector} proves uniqueness inductively because
the auxiliary constant acts as the displayed nonzero scalar on that sector.
Equivalently its sector inverse has identity
$(1+\eta\eta'\eta_2\eta_3)/2$, rather than the identity of the full algebra.
The auxiliary higher coefficients preserve the corresponding sector, so each
exact remainder stays there.

The implementation orders monomials by total degree, stores positive-degree
auxiliary terms in that order, and updates future remainders immediately after
solving a coefficient. Updates stop at the chosen total or componentwise
cutoff and never revisit a solved
coefficient. Only physical triples and their eight parity slots are stored.

\section{Precision, timings, and PBW comparison}
\label{sec:timings}

Machine arithmetic uses \texttt{std::complex<double>}; higher precision uses
an owning MPC type with MPFR/GMP and the same templated mathematical routines.
At 40 digits it uses 136 bits. This does not mean 40 certified output digits:
the mode-action solves still use double-precision column-rank selection and
LU preconditioning. Outer branching instead uses the stored recursion of
Section~\ref{sec:outer-recursion}, with local arithmetic at working precision.
At decimal precision $D$, the selected-system residual target is
$10^{-\max(20,D-15)}$ and the full-system residual guard is
$10^{-\max(15,D-20)}$. Sparse pruning uses $10^{-\max(20,D-20)}$.
Native action residuals are required below $10^{-8}$, and native sparse
pruning uses $10^{-13}$. Each outer Ward residual is the absolute value of the sum of its
terms divided by one plus the sum of their absolute values; its guard is
$10^{-8}$ at machine precision and $10^{-\max(15,D-20)}$ at $D$ digits.
At most 30 refinement iterations are attempted. Ill-conditioned rank selection
or failed convergence raises an error rather than silently changing the basis.

The sector residual is evaluated during recovery, before vacuum restoration,
on the coefficients of $\FF/F_{\rm vac}^2$. It is the distance of a parity vector from
$[v+\eta\eta'(\eta_2\eta_3)\st v]/2$, divided by the largest of one and
its component magnitudes. The default policy stops above $10^{-8}$.
It is not recomputed on the final restored arrays.
Ordinary default recovery projects accepted small residuals; inserted default
recovery retains its coefficients. The explicit \texttt{record} policy retains
unprojected coefficients in both modes and records the residual. Timing runs
below use this policy, so completion alone is not an accuracy certificate.

\subsection{Fresh independent-edge level-10 measurements}
Both production pipelines use the stored outer recurrence. The parameters are
\begin{equation}
 b=7/5,\quad(P_1,P_2,P_3)=(11/23,13/29,17/31),\quad p_1=f=0,
 \quad(\eta,\eta')=(+,+)\text{ or }(+,-).
\end{equation}
Each sector and precision is a separate process with empty numerical caches.
Runs are sequential; within-process reuse remains enabled. The build uses
Apple clang 17 with \texttt{-O3} on macOS arm64 with Accelerate.
Fermion Ward sums and Schottky series use exact GMP rationals.
The internal timer includes branching, CCY, products, assembly, direct fermion,
convolution, and Schottky restoration. A parent timer additionally includes
process startup, final output serialization, and exit. Compilation is excluded.
The executable and source hashes, commands, diagnostics, stage times, and full
physical coefficient arrays are saved under
\path{C++/results/per_edge_level10_2026-09-11/optimized_full_pipeline_40dps/}.

% Generated by C++/tools/render_per_edge_notes.py from completed run records.
These runs impose $N=10$ independently on each edge, retaining all
$2{,}541$ monomials and $20{,}328$ parity components in each sector.
They use 40 decimal digits of working precision (136 bits), with the
fermion Ward sums and Schottky series evaluated exactly. All table entries are seconds.
\begin{center}\small
\begin{tabular}{@{}lrr@{}}
\toprule
Stage & $\eta\eta'=1$ & $\eta\eta'=-1$\\
\midrule
Initial mode actions & 9.497 & 9.594\\
Outer branching recurrence & 0.441 & 0.450\\
Middle recurrence & 0.000 & 0.013\\
Reduced CCY & 296.593 & 1460.243\\
Virasoro products & 26.759 & 88.559\\
Branch assembly & 0.356 & 0.577\\
Direct fermion & 5.220 & 5.221\\
Sector division & 3.115 & 3.130\\
Schottky vacuum and square & 72.379 & 72.492\\
Vacuum restoration & 0.375 & 0.380\\
Other pipeline overhead & 1.801 & 14.620\\
Internal total & 416.537 & 1655.279\\
Startup, serialization and exit & 0.081 & 0.081\\
\midrule
Full wall time & \textbf{416.619} & \textbf{1655.360}\\
\bottomrule
\end{tabular}
\end{center}
The full wall times are \textbf{6 min 57 s} and \textbf{27 min 35 s}.
Other pipeline overhead is the difference between the independently timed
internal total and the named stages; it includes work and cache cleanup
outside their timers. Branching, including initial actions but excluding
the middle recurrence, takes $9.938$ s and $10.044$ s. The outer Ward
recurrence itself takes $0.441$ s and $0.450$ s.
\begin{center}\small
\begin{tabular}{@{}lrr@{}}
\toprule
Recorded work & $\eta\eta'=1$ & $\eta\eta'=-1$\\
\midrule
Direct boundary values & 96 & 192\\
Stored recursive branching values & 1,704 & 3,408\\
Assembled branch tuples & 899 & 1,591\\
Virasoro blocks & 1,798 & 1,620\\
Reused transposed products & 0 & 781\\
Global-seed summands & 42,624,112 & 1,740,820,224\\
CCY transitions & 51,242,176 & 839,918,556\\
\bottomrule
\end{tabular}
\end{center}
The ordinary wall time was $508.004$ s before the latest CCY reuse changes;
its CCY stage decreased from $390.230$ s to $296.593$ s. The former
inserted run was interrupted and had projected to about $4.5$ hours.
That forecast is not a completed baseline measurement. Interrupted and
overlapping attempts are excluded from the tables above.
\subsubsection*{Directed checks and accuracy limits}
Both saved outputs contain the complete requested box and only finite
physical components. Source and executable hashes match the launch records.
The maximum recorded recovery-sector residuals are $2.226\times10^{-15}$
and $8.837\times10^{-16}$, respectively. These test sector consistency
before vacuum restoration, not independent physical-coefficient accuracy.

Comparing all $20{,}328$ ordinary components with the saved pre-change
40-digit result gives maximum scaled difference $2.107\times10^{-20}$
and absolute difference $8.907\times10^{-9}$; coefficients reach
$4.166\times10^{12}$. Here the scale is $\max(1,|a|,|b|)$.
This checks the implementation change at the same working precision.
The full inserted box has no completed pre-change or independent reference.

At independent cutoff 2, both physical pipelines were compared with saved
results on all 45 monomials; after the CCY optimization their scaled
differences from the saved pre-change box are below $2.97\times10^{-35}$.
Individual inserted Virasoro factors at edge cutoffs 6 and 8 agree with
the former implementation within $8.88\times10^{-32}$ in scaled difference.
Exact checks compare small Schottky and fermion boxes with projections of
their total-degree series and check the shifted-diagonal product on an
asymmetric domain. No additional physical PBW block was computed.

\begingroup\raggedright
The directed optimization evidence is saved under
\path{C++/results/ccy_vertex_reuse_2026-09-11/}; output inspection and
cutoff validation are in \path{C++/results/per_edge_level10_2026-09-11/}.
\par\endgroup


\clearpage
\subsection{Measured C++ comparison at independent level 5}
\label{sec:pbw-measured}
% Generated by C++/tools/render_optimized_pbw_notes.py from fresh C++ measurements.
Both numerical implementations are C++17 with MPC at 40-digit working
precision (136 bits). At the same parameters, fresh runs with
$q_1,q_2,q_3$ each cut off at level 5 give the following full wall times.
Each result contains 396 monomials and 3,168 parity components, including
$q_1^5q_2^5q_3^5$. Double Virasoro includes all stages of physical-block
recovery. Processes are separate and sequential; numerical caches start
empty, within-run reuse is enabled, and compilation is excluded.
Python only launches the executables and compares saved outputs.
\begin{center}\small
\begin{tabular}{@{}lrrr@{}}
\toprule
Sector & C++ PBW (s) & C++ double Virasoro (s) & PBW/DV\\
\midrule
$\eta\eta'=1$ & 4.77 & 3.10 & $1.54\times$\\
$\eta\eta'=-1$ & 8.61 & 7.32 & $1.18\times$\\
\bottomrule
\end{tabular}
\end{center}
The measured double-Virasoro advantage is modest at this cutoff:
$1.54\times$ and $1.18\times$ in the two sectors.
These are single-run timings, not statistical performance bounds.

\paragraph{Direct C++ PBW implementation.}
The physical three-point tensors are evaluated by SCA Ward recursion and
contracted with inverse physical Gram matrices. The negative-sector
command uses the label \texttt{inserted} to match the other executable;
this direct PBW calculation itself contains no auxiliary insertion.
Descendant words are interned once and states use integer labels.
The implementation caches canonical words, levels, parities, binomial
coefficients and mode actions, sharing action data between the two vertex
signs. Ward values use three-integer keys. Each symmetric bilinear Gram
matrix is built from one triangle and its inverse is retained. When
$\eta=\eta'$, the second vertex reuses the first tensor. Successive
axis contractions index the tensors directly, with reusable MPC scratch
and no stored permutation tables. Output is serialized once.

Gram LU elimination and substitution use working precision throughout;
only pivot comparisons convert magnitudes to double. There is no
double-precision LU preconditioner in this direct PBW implementation.
The numerical tensors use bilinear, not Hermitian, contraction.
\begin{center}\small
\begin{tabular}{@{}lrr@{}}
\toprule
C++ PBW stage (s) & $\eta\eta'=1$ & $\eta\eta'=-1$\\
\midrule
Gram entries & 0.009 & 0.009\\
Gram inversion & 0.007 & 0.006\\
Vertex tensors and Ward recursion & 3.666 & 7.413\\
Tensor contractions & 0.999 & 0.999\\
Internal total & 4.694 & 8.454\\
\bottomrule
\end{tabular}
\end{center}
Vertex construction accounts for 78.1\% and 87.7\% of internal time. Peak resident memory is 102.9 and 191.3 MB (decimal).

The double-Virasoro benchmark includes the two latest CCY changes:
grouping inserted assembly by recursion shift and middle descendant
levels to reuse the amplitude times middle vertex; and dense complete-vertex
caches, with reusable MPC scratch during propagation. These evaluate
the same sum in \eqref{eq:forward-assemble}. Ordinary three-edge global
assembly retains its production formula. This isolated benchmark build
is separate from the executable used for the measured level-10 table above.

\paragraph{Coefficient comparison.}
All 3,168 components per sector were compared. For C++ PBW versus
C++ double Virasoro, the maximum scaled differences,
$|a-b|/\max(1,|a|,|b|)$, are $8.598\times10^{-21}$
and $1.395\times10^{-20}$; the absolute maxima are
$3.951\times10^{-16}$ and $1.724\times10^{-16}$.
This is agreement at roughly 20 digits on this scale, not 40 certified
output digits. To check the port separately, comparison with the saved
Python PBW coefficients gives maximum scaled differences $2.817\times10^{-38}$
and $1.732\times10^{-36}$. No Python block was rerun.

\begingroup\raggedright
The direct code is \path{C++/include/ramond/direct_pbw.hpp}, with driver
\path{C++/drivers/pbw_main.cpp}. Timings, coefficients, comparison records
and source/build hashes are under
\path{C++/results/direct_pbw_cpp_2026-09-12/}.
\par\endgroup


\subsection{C++ PBW estimate at independent level 10}
\label{sec:pbw-estimate}
The updated estimates use the measured 136-bit stages of the optimized
C++/MPC physical SCA calculation in Section~\ref{sec:pbw-measured}.
No independent-edge level-10 PBW block is evaluated.
The dimensions used for work counts are
\begin{equation}
 \sum_{a\ge0}d_{\NS}(a)x^a=\prod_{m\ge1}\frac{1+x^{2m-1}}{1-x^{2m}},\qquad
 \sum_{l\ge0}d_{\R}(l)y^l=\prod_{m\ge1}\frac{1+y^m}{1-y^m}.
\end{equation}
The NS index is twice the level and the Ramond dimension is per parity.
Counts of required tensor entries, reused Gram entries, and contractions are
computed from these generating functions and the actual PBW loop bounds.
The sign choice does not change basis sizes. For $\eta=\eta'$ the optimized
implementation constructs one vertex tensor and reuses it at the second
vertex, halving the requested form entries. Gram and contraction work is
unchanged. Each stage is scaled by its own work count, as detailed below;
the two time scenarios describe model sensitivity, not a confidence interval
or a guaranteed runtime bound.

% Generated by C++/tools/render_per_edge_notes.py from saved count and calibration records.
\begin{samepage}
For the independent-edge level-10 box, the required PBW work is
\begin{center}\small
\begin{tabular}{@{}lr@{}}
\toprule
Quantity & Count\\
\midrule
Physical monomials & 2,541\\
Largest NS Gram dimension & 161\\
Largest R Gram dimension, per parity & 232\\
Entries in distinct Gram matrices & 450,469\\
Cubic Gram sum (inversion-cost proxy) & 78,765,711\\
Requested vertex entries, negative sector & 1,318,075,412\\
Dense contraction multiplications & 254,809,434,024\\
Largest single vertex tensor & 8,665,664 entries\\
\bottomrule
\end{tabular}
\end{center}
\end{samepage}
The per-edge Gram dimensions are unchanged from total level 10, but the
box includes many more combinations of high levels. For dimensions
$d_1,d_2,d_3$, two vertices and the two parity sectors request
$4d_1d_2d_3$ form entries; successive Gram contractions and the final
pairing use $2d_1d_2d_3(d_1+d_2+d_3+1)$ complex multiplications.
The sum of dimension cubes is a dense-inversion work proxy, not an exact
operation count. Identical vertices in the positive sector can share
form evaluations; only the fitted form-entry cost is allowed to fall
by a factor between one half and one. Gram and contraction costs remain.
\begin{center}\small
\begin{tabular}{@{}lcc@{}}
\toprule
136-bit PBW model & $\eta\eta'=1$ & $\eta\eta'=-1$\\
\midrule
Lean & $28.8$--$63.4$ h & $57.6$--$63.4$ h\\
Full & $97.3$--$196.0$ h & $119.3$--$196.0$ h\\
\bottomrule
\end{tabular}
\end{center}
These are estimates for the existing Python/FLINT PBW implementation.
Saved 384-bit PBW timing fits are rescaled using synthetic scalar and
matrix operations in its actual arithmetic wrappers at 136 and 384 bits.
The arithmetic-only calibration took $1.30$ s and evaluated no physical
block, Gram matrix or Ward identity. Measured 136/384-bit time ratios
were about $0.91$--$0.92$ for scalar operations and $0.51$--$0.73$ for
matrix products. The estimates allow unchanged Python/cache overhead.

The lean fit omits explicit cubic-Gram and contraction features; the full
fit includes them. Their spread exposes model sensitivity, rather than
forming a confidence interval. Larger matrix shapes, internal Ward work
and memory pressure may change these costs; sufficient memory is assumed.
The C++ and PBW comparison now uses the same 136-bit working precision,
but different implementations (MPC versus Python/FLINT). No full 136-bit
PBW block was timed for this box.

\begingroup\raggedright
Integer counts, saved-fit provenance and the arithmetic calibration are in
\path{C++/results/per_edge_level10_2026-09-11/work_estimates.json} and
\path{C++/results/per_edge_level10_2026-09-11/pbw_40digit_arithmetic_calibration.json}.
\par\endgroup


\subsection{Reproduction and evidence}
Under \path{C++/include/ramond/}, the stages are separated into
\path{free_field.hpp} and \path{branching.hpp} (actions
and branching), \path{ccy.hpp} (reduced Virasoro factors),
\path{schottky.hpp} (vacuum), \path{fermion.hpp} (direct auxiliary sewing),
and \path{pipeline.hpp} (branch-specific products, assembly, and recovery).
The separate \path{direct_pbw.hpp} and \path{C++/drivers/pbw_main.cpp}
implement direct physical SCA sewing for the comparison.
Partition lists, solved action data, outer primary coefficients, and auxiliary
Ward values are reused within a run. Pole, fusion, inverse-denominator,
global-core and complete inserted-vertex caches belong
to their individual Virasoro engine; distinct branching weights do not share
those numerical values. This distinction fixes the meaning of fresh runs
and reuse in the timing tables.

\begin{samepage}
From the repository root:
\begin{verbatim}
make -C 'C++' -j2
python3 'C++/tools/time_current_pipelines.py' --levels 10 --dps 40 \
  --truncation per-edge --output '/tmp/ramond_per_edge_level10'
\end{verbatim}
\end{samepage}
The coefficient convention in JSON is $[a,l,l,d]$ for
$q_1^{a/2}q_2^lq_3^{d/2}$, with even $d$. Under the independent-edge
cutoff, $0\le a,d\le2N$ and $0\le l\le N$; under a total cutoff,
$a+2l+d\le2N$. The JSON identifies the selected truncation and records
the independent limits as \texttt{q\_level\_cutoffs}.
Parity slot is $\epsilon_1+2\epsilon_2+4\epsilon_3$.
Directed checks of the stored recurrence cover all 768 ordinary-support and
880 inserted-support outer coefficients through level 5. Agreement with the
saved reference coefficients is within $3.76\times10^{-25}$ in scaled
difference at 40 digits, including equal Ramond momenta. The equal-momentum
case exercises vanishing first prefactors; each low-level run checks 608
second Ward identities across both vertex signs and Ramond parities.
The independent-edge output checks and their limits are reported above.
The earlier production checks used stored physical references. The fresh
independent-edge level-5 PBW comparison is recorded in
Section~\ref{sec:pbw-measured}; no level-10 PBW computation was used for the
updated estimates. The C++ PBW benchmark and its count-only estimate are
reproduced by
\begin{verbatim}
make -C 'C++' bin/pbw
python3 'C++/tools/benchmark_direct_pbw.py'
python3 'C++/tools/estimate_direct_pbw.py'
\end{verbatim}
The benchmark uses the saved isolated double-Virasoro executable described
in Section~\ref{sec:pbw-measured}; its build recipe is
\path{C++/experiments/pbw_vs_double_virasoro_level5_2026-09-12/build_ccy.py}.
Regenerate all timing and PBW tables from saved records,
without evaluating blocks, using
\begin{verbatim}
python3 'C++/tools/render_per_edge_notes.py'
\end{verbatim}

\begin{thebibliography}{9}
\bibitem{CCY} M. Cho, S. Collier and X. Yin,
\emph{Recursive Representations of Arbitrary Virasoro Conformal Blocks},
\href{https://arxiv.org/abs/1703.09805}{arXiv:1703.09805}.
\end{thebibliography}
\end{document}
